\chapter{Snapshot testing}
\label{ch:snapshot}

\begin{aside}
Render. Compare to a recorded screen. Diff. That's a snapshot
test --- the most efficient way to test UI for ``it still looks
right''.
\end{aside}

\section{The idea}

A snapshot test:

\begin{enumerate}[leftmargin=*]
  \item Construct a known state.
  \item Render to an in-memory canvas (string backend).
  \item Compare against a stored ``golden'' file.
  \item Pass if equal; fail with a diff otherwise.
\end{enumerate}

When the UI changes intentionally, regenerate the goldens.
When a refactor unintentionally changes pixels, the diff is
obvious in code review.

\section{Capturing}

\begin{lstlisting}
TEST_CASE("button renders") {
    Context ctx;
    auto el = button("Click me");
    auto snap = render_to_string(el, 20, 3);
    EXPECT_SNAPSHOT(snap, "button.snap");
}
\end{lstlisting}

The first run creates \code{button.snap}; subsequent runs
compare.

\section{Format}

A snapshot is a plain-text grid: each row a line, columns
separated by nothing. Styles are encoded as inline markers
(or omitted, for layout-only snapshots).

For style verification, an extended format includes per-cell
foreground/background/attribute lines. \maya{}'s snapshot
helper supports both.

\section{Diff display}

When a snapshot fails, the diff is shown side-by-side or
with red/green highlighting per character. The reader can
spot a one-cell shift instantly.

\section{Regenerating goldens}

\begin{lstlisting}
UPDATE_SNAPSHOTS=1 ctest
\end{lstlisting}

Sets every snapshot to the current output. Review the diff
before committing.

\section{Snapshot stability}

Snapshots break on:

\begin{itemize}[leftmargin=*]
  \item Font metric changes (cell-pixel ratio differs across
    machines).
  \item Timestamp inclusion (don't put dates in snapshots).
  \item Random data (seed your RNG).
  \item Locale-dependent formatting (force a locale).
\end{itemize}

\maya{}'s snapshot helper provides hooks to mask volatile
regions:

\begin{lstlisting}
snap.mask({0, 0, 8, 1}, "TIME");
\end{lstlisting}

The masked region renders as the literal string for
comparison.

\section{Layout-only snapshots}

For tests that only care about structure, snapshot the
element tree as a text outline:

\begin{lstlisting}
auto tree = dump_tree(el);
EXPECT_SNAPSHOT(tree, "button.tree");
\end{lstlisting}

The output is something like:

\begin{lstlisting}
button {
  text("Click me")
  width: 20
  height: 1
}
\end{lstlisting}

Style and layout are tested separately.

\section{Property-based snapshots}

Combine snapshots with property-based testing: generate
random states, render, ensure no crash, snapshot canonical
states.

\maya{} ships a \code{property::run} helper that integrates
with rapidcheck.

\section{Visual snapshot}

For terminals that support image protocols (Kitty), render to
a PNG via a headless image protocol and store that. Compare
PNGs pixel-by-pixel. Catches font-rendering regressions but
needs identical font on every CI machine.

Most teams stop at the text-grid snapshot.

\section{Snapshot review in code review}

Treat \code{.snap} files as code: review them. A change that
``just'' updates a snapshot is often the right thing to
review most carefully --- the snapshot is the contract.

\section{Recap}

\begin{recap}
\begin{itemize}[leftmargin=*]
  \item Render + compare against stored golden = snapshot.
  \item Update with \code{UPDATE\_SNAPSHOTS=1}.
  \item Mask volatile regions.
  \item Layout-only snapshots for structural tests.
  \item Review snapshot diffs as you would code.
\end{itemize}
\end{recap}

\section*{Workshop}

\begin{enumerate}[leftmargin=*]
  \item Add a snapshot test for your app's main screen.
    Commit the golden.
  \item Make a visual tweak (border style). Watch the
    snapshot diff and update.
  \item Add masks for the timestamp in your status bar.
\end{enumerate}
